% !TEX root = ../bb_AiRa_Kappaleet.tex

% ö = \%c3\%b6
% ä = \%C3\%A4

\xdef\sectiontitle{\IIIsectiontitle} %aiheuttama
\TOCrefs

%%%%%%%%%%%%%%%
\begin{frame}[label=3_3_a]%[label=3_3_TOC1]
\frametitle{\texorpdfstring{\culine[\sectiontitleulinecolor]{\sectiontitle}}{\sectiontitle} }
\label{\TOCname}

\vspace{-0.5cm}\label{\TOCname} % vertical spacing is off, 0.7cm doesn't work

\begin{itemize}[leftmargin=0.5cm,itemsep=\chapterTOCitemsep]
    \item[3.1]<1-> \hyperlink{3_1_a}{\IIIaSubsectiontitle}
    \item[3.2]<1-> \hyperlink{3_2_a}{\IIIbSubsectiontitle}	
    \item[3.3]<1-> \hyperlink{3_3_a}{\CDAlert[blue]{\IIIcSubsectiontitle}}	
    \item[3.4]<1-> \hyperlink{3_4_a}{\IIIdSubsectiontitle}
    \item[3.5]<1-> \hyperlink{3_5_a}{\IIIeSubsectiontitle}
    \item[3.6]<1-> \hyperlink{3_6_a}{\IIIfSubsectiontitle}
    \item[3.7]<1-> \hyperlink{3_7_a}{\IIIgSubsectiontitle}
\end{itemize}

%\endofslide 
\end{frame}
%%%%%%%%%%%%%%%

\xdef\sectiontitle{\IIIcSubsectiontitle} 


%%%%%%%%%%%%%%%
\begin{frame}
\frametitle{\texorpdfstring{\culine[\sectiontitleulinecolor]{\sectiontitle}}{\sectiontitle} }
\framesubtitle{Ytimen muoto}%Shape of the nucleus
\appear{}

\begin{itemize}
	\item Valtaosa ytimistä on \CDAlert[red]{lähes pallomaisia}
	\item Lähinnä vain harvinaisten maametallien ytimet \appear{$(Z=57-71)$} \appear{ovat \CDAlert[blue]{ellipsoideja}}\appear{, joiden pääakselin pituus voi vaihdella n.\ 20\% verran sivuakselin pituudesta}

\item Näiden painavien ytimien  sisäelektronien aaltofunktiot jopa vähän ulottuvat ytimen sisään. \appear{Poikkeamat pallomaisesta muodosta}\appear{ vastaavat \Alert[fuchsia]{ poikkeamia ydinvarausten tiheysjakaumissa}}\appear{, ja näkyvät pieninä eroina ydinten energiatasoissa}

\item \CDAlert[blue]{Ytimen muoto} näkyy myös nollasta poikkeavana keskimääräisenä   \CDAlert[tuniblue]{sähköisenä kvadrupolimomenttina}:
\[
< Q >  \equal \appear{Z} \appear{\nint} \appear{\psi^{*}} \appear{\textrm{[[} 3} \appear{\textrm{((}z_{a v}^{2} \textrm{))}_{a v}}  \minus  \appear{\textrm{((}x^{2}+y^{2}+z^{2} \textrm{))}_{a v}\textrm{]]}} \appear{\psi}  \, \appear{d V}
\]
\end{itemize}

%\xdef\loc{south east}
%\begin{tikzpicture}[remember picture,overlay] % anchor=south
%    \node (A) [yshift=-0.25*\figYshift,xshift=\figXshift, anchor=\loc] at (current page.\loc) {\includegraphics[width=7cm]{Figures_booklet/SOM_12_Nuclear_physics_figures/SOM_Nuclear_physics_Figure_1.png}}; % 8cm max height
%\end{tikzpicture}

\endofslide 
%\endofsection
\end{frame}
%%%%%%%%%%%%%%%

%%%%%%%%%%%%%%%
\begin{frame}
\frametitle{\texorpdfstring{\culine[\sectiontitleulinecolor]{\sectiontitle}}{\sectiontitle} }
\framesubtitle{Ytimen muoto}%Shape of the nucleus
\appear{}
%\vspace{2.5mm}
\begin{minipage}{10.0cm}
\begin{itemize}
	\item Jos kvadrupolimomentin arvo on \CDAlert[blue]{positiivinen}\appear{, ytimen muoto on  \CDAlert[blue]{pitkänomainen ellipsoidi}}\appear{, kuin pystyssä oleva vesimeloni}
	\appear{\xdef\loc{south east}
\begin{tikzpicture}[remember picture,overlay] % anchor=south
    \node (A) [yshift=-0.15*\figYshift,xshift=\figXshift, anchor=\loc] at (current page.\loc) {\includegraphics[height=8.75cm]{Figures_booklet/SOM_12_Nuclear_physics_figures/SOM_Tipler_Nuclear_physics_figure_10.png}}; % 8cm max height
\end{tikzpicture}
\xdef\loc{south east}
\begin{tikzpicture}[remember picture,overlay] % anchor=south
    \node (A) [yshift=-0.15*\figYshift,xshift=-3.5cm, anchor=\loc] at (current page.\loc) {\includegraphics[width=4cm]{Figures_booklet/SOM_12_Nuclear_physics_figures/SOM_Tipler_Nuclear_physics_figure_11_cropped.png}}; % 8cm max height
\end{tikzpicture}\CR[ref = t, pos = left]{}} 
	\item Jos momentti on  \CDAlert[red]{nolla}\appear{ ydin on \CDAlert[red]{pallomainen}} 
%	\item If the quadrupole moment is \CDAlert[fuchsia]{negative}\appear{, then the shape is \CDAlert[fuchsia]{oblate ellipsoid}, i.e. more like a 'flattened pumpkin'}
\end{itemize}
\end{minipage}
\vspace{2.5mm}

\begin{minipage}{6.5cm}
\begin{itemize}
	\item Jos momentti on \CDAlert[fuchsia]{negatiivinen}\appear{, niin ytimen muoto on \CDAlert[fuchsia]{(navoiltaan)}\\ \CDAlert[fuchsia]{litistynyt ellipsoidi}}
\end{itemize}
\end{minipage}
%\vspace{2.5mm}



\endofslide 
%\endofsection
\end{frame}
%%%%%%%%%%%%%%%

%%%%%%%%%%%%%%%
\begin{frame}
\frametitle{\texorpdfstring{\culine[\sectiontitleulinecolor]{\sectiontitle}}{\sectiontitle} }
\framesubtitle{Ydinmallit}%Nuclear models
\appear{}

%\seminormal
\begin{itemize}[itemsep=1.7mm]
	\item Atomin ydintä voidaan kuvata usealla teoreettisella mallilla. \appear{Kaksi eniten käytettyä ovat:}
\item[] \begin{itemize}
	\item \CDAlert[red]{Nestepisaramalli}%Liquid drop model
\begin{itemize}
	\item Olettaa ytimen olevan (kokoonpuristumaton) pyöreä nestepisara
\item Mahdollistaa \CDAlert[tunired]{ydinten sidosenergioiden laskemisen}
\end{itemize}

\item \CDAlert[blue]{Kuorimalli}
\item<.->[] \begin{itemize}
	\item Soveltaa 'elektroni kvanttikaivossa'-mallia ytimen tapaukseen
\item Ennustaa esim.\ vakaimman nuklidin $A$-arvon\appear{, ja pohjautuu ns.\ \CDAlert[tuniblue]{taikalukuihin}}
\end{itemize}

\end{itemize}

\item Molemmilla malleilla on vahvuutensa. \appear{Vaikka ne ovat osittain keskenään yhteensopivia}\appear{, ne ovat vain \CDAlert[fuchsia]{approksimaatioita}}\appear{, ja perustuvat tiettyjen oletusten ja yksinkertaistuksien varaan. } \appear{Vaihtoehtoisille malleille on vielä tilaa/kysyntää atomin ytimen ymmärtämiseksi}
\end{itemize}

%\xdef\loc{south east}
%\begin{tikzpicture}[remember picture,overlay] % anchor=south
%    \node (A) [yshift=-0.25*\figYshift,xshift=\figXshift, anchor=\loc] at (current page.\loc) {\includegraphics[width=7cm]{Figures_booklet/SOM_12_Nuclear_physics_figures/SOM_Nuclear_physics_Figure_1.png}}; % 8cm max height
%\end{tikzpicture}

\endofslide 
%\endofsection
\end{frame}
%%%%%%%%%%%%%%%

%%%%%%%%%%%%%%%
\begin{frame}
\frametitle{\texorpdfstring{\culine[\sectiontitleulinecolor]{\sectiontitle}}{\sectiontitle} }
\framesubtitle{Nestepisaramalli}%The liquid drop model
\appear{}

%
\begin{itemize}
	\item Koska ytimen tiheys oli pikemminkin radiaalinen kuin täysin homogeeninen\appear{, ja koska kaikkien ytimien tiheys on likimain vakio}\appear{, on luonnollista kuvata ydintä \Alert[red]{nestepisaramallin} avulla}

\item Analogia nestepisaraan toimii suhteellisen hyvin: 
\item[] \begin{itemize}
	\item Neste koostuu tiiviisti pakkautuneista\appear{, järjestäytymättömistä molekyyleistä}\appear{, jotka vetävät \CDAlert[red]{toisiaan puoleensa} suuremmilla etäisyyksillä} 
	\item Pienemmillä etäisyyksillä ne \CDAlert[red]{hylkivät toisiaan}\appear{, johtaen \CDAlert[red]{puristumattomuuteen}, kuten esim.\ veden tapauksessa}
\end{itemize}


\item Jotta nestemolekyyli voidaan siirtää pisaran pinnalle\appear{, molekyyli täytyy vetää erilleen}\appear{ muista sitä puoleensavetävistä molekyyleistä.} \appear{Tämä vaatii energiaa}\appear{, jolloin pallomainen muoto} \appear{johtaa matalimman energian tilaan}\appear{, ts.\ minimoi pinta-alan}
\end{itemize}

%\xdef\loc{south east}
%\begin{tikzpicture}[remember picture,overlay] % anchor=south
%    \node (A) [yshift=-0.25*\figYshift,xshift=\figXshift, anchor=\loc] at (current page.\loc) {\includegraphics[width=7cm]{Figures_booklet/SOM_12_Nuclear_physics_figures/SOM_Nuclear_physics_Figure_1.png}}; % 8cm max height
%\end{tikzpicture}

\endofslide 
%\endofsection
\end{frame}
%%%%%%%%%%%%%%%

%%%%%%%%%%%%%%%
\begin{frame}
\frametitle{\texorpdfstring{\culine[\sectiontitleulinecolor]{\sectiontitle}}{\sectiontitle} }
\framesubtitle{Nestepisaramalli}%The liquid drop model
\appear{}

\semismall
\begin{itemize}
	\item Pisaran sisällä olevia hiukkaset vetävät puoleensa kaikkia läheisyydessä olevia hiukkasia

\item Jos kaikki hiukkaset olisivat pisaran sisällä\appear{, niin sidosenergia olisi suoraan verrannollinen hiukkasten lukumäärään.}%
\appear{\xdef\loc{south east}%
\begin{tikzpicture}[remember picture,overlay] % anchor=south
    \node (A) [yshift=-0.25*\figYshift,xshift=\figXshift, anchor=\loc] at (current page.\loc) {\includegraphics[width=5cm]{Figures_booklet/SOM_12_Nuclear_physics_figures/Tikz_SOM_Nuclear_liquid_drop.pdf}};% 8cm max height
\end{tikzpicture}}
 \appear{\Alert[fuchsia]{Pinnalla olevilla hiukkasilla on kuitenkin vähemmän naapureita}}\appear{, eli vähemmän sidoksia toisten hiukkasten kanssa} \implies{ Sidosenergia pienenee }

\end{itemize}

% 6.5 cm = midway, 10 cm = 3/4 
\vspace{2.5mm}
\begin{minipage}{8.5cm}
\begin{itemize}
	\item Muodostetaan seuraavaksi \CDAlert[fuchsia]{yleinen kaava} nestepisaran tavoin käyttäytyvän ytimen \CDAlert[fuchsia]{sidosenergialle} 

\item Yhteensä löytyy \CDAlert[red]{viisi vaikuttavaa tekijää} \appear{(eli viisi termiä)} \appear{jotka vaikuttavat sidosenergian suuruuteen}

\item Määritellään vakiot $c_{1}\appear{, c_{2}}\appear{, c_{3}, c_{4}$, ja $c_{5}$} \appear{sovittamalla kokeellinen data muodostettuun malliin [jonka on oltava $f(A,N,Z)$]}

\item Tavoitteena on rakentaa \CDAlert[blue]{semi-empiirinen kaava} tunnettujen ytimien sidosenergioille
\end{itemize}
\end{minipage}
%\vspace{2.5mm}




\endofslide 
%\endofsection
\end{frame}
%%%%%%%%%%%%%%%

%%%%%%%%%%%%%%%
\begin{frame}
\frametitle{\texorpdfstring{\culine[\sectiontitleulinecolor]{\sectiontitle}}{\sectiontitle} }
\framesubtitle{Tilavuustermi}%Volume term
\appear{}

%
\begin{itemize}
	\item \CDAlert[red]{Kokonaan pisaran sisällä olevat hiukkaset}\appear{ ovat täysin naapureidensa ympäröimiä.} \appear{ Jos kaikki hiukkaset olisivat pisaran sisällä}\appear{, niin sidosenergia olisi suoraan verrannollinen hiukkasten} \appear{(eli ytimen nukleonien)} \appear{lukumäärään jotka muodostavat pisaran.}
\[
BE \propto \appear{c_{1}} \appear{A}
\]

\end{itemize}

% 6.5 cm = midway, 10 cm = 3/4 
\vspace{2.5mm}
\begin{minipage}{8.5cm}
\begin{itemize}
\item Verrannollisuuskerroin $c_{1}$ voitaisiin teoreet\-ti\-sesti laskea\appear{, mutta on helpompaa vain muodostaa malli ja lopuksi \Alert[fuchsia]{sovittaa  vakiot kokeelliseen dataa}}

	\item Eli \CDAlert[red]{tilavuustermi} on muotoa  \[
c_{1} A
\]
\end{itemize}
\end{minipage}
%\vspace{2.5mm}


\xdef\loc{south east}%
\begin{tikzpicture}[remember picture,overlay] % anchor=south
    \node (A) [yshift=-0.25*\figYshift,xshift=\figXshift, anchor=\loc] at (current page.\loc) {\includegraphics[width=5cm]{Figures_booklet/SOM_12_Nuclear_physics_figures/Tikz_SOM_Nuclear_liquid_drop_interior.pdf}};% 8cm max height
\end{tikzpicture}

\endofslide 
%\endofsection
\end{frame}
%%%%%%%%%%%%%%%

%%%%%%%%%%%%%%%
\begin{frame}
\frametitle{\texorpdfstring{\culine[\sectiontitleulinecolor]{\sectiontitle}}{\sectiontitle} }
\framesubtitle{Tilavuustermi}%Volume term
\appear{}

\semismall
\begin{itemize}[itemsep=1.5mm]
	\item K: Onko sovitus/verrannollisuukertoimiin $c_{i}$ pohjautuva empiirinen malli \appear{hyödyllinen?}
\item V: \CDAlert[fuchsia]{Malli kuvaa monimutkaista systeemiä}\appear{, ja se perustuu mittaustuloksiin}
\appear{\xdef\loc{south east}
\begin{tikzpicture}[remember picture,overlay] % anchor=south
    \node (A) [yshift=-0.35*\figYshift,xshift=\figXshift, anchor=\loc] at (current page.\loc) {\includegraphics[width=7cm]{Figures_booklet/SOM_12_Nuclear_physics_figures/SOM_Nuclear_physics_Figure_14.png}}; % 8cm max height
\end{tikzpicture}\CR[ref = h]{}}
\item Sidososuudet muodostavat monimutkaisen funktion $\mathrm{N}$:n \appear{ja $\mathrm{Z}$:n suhteen}\appear{, sen matemaattista muotoa ei voida eksplisiittisesti nähdä}

	\item Yksinkertaisen $Z$:sta ja $N$:sta riippuvan polynomisovitteen
\[
c_{1} \appear{Z^{2}}  \plus \appear{c_{2} N}  \plus \appear{c_{3}} %\hspace{4cm}
\]
\end{itemize}
% 6.5 cm = midway, 10 cm = 3/4 
%\vspace{2.5mm}
\begin{minipage}{6.5cm}
\begin{itemize}[itemsep=1.5mm]
\item[] avulla ei saavuteta hyvää ennus\-tet\-ta\-vuutta\appear{, tai saada hyvää fysikaalista in\-tui\-tio\-ta ytimen käyttäytymiseen}

\item Koska nestepisaramalli sopii niin hyvin kokeellisiin tuloksiin\appear{, on järkevää olettaa että malli ja sen alkuoletukset ovat toimivia}

\item Mallin avulla voidaan myös \CDAlert[fuchsia]{ennustaa} tuntemattomien nuklidien sidosenergioita \appear{(ennustettavuus)}
\end{itemize}
\end{minipage}
%\vspace{2.5mm}

\endofslide 
%\endofsection
\end{frame}
%%%%%%%%%%%%%%%

%%%%%%%%%%%%%%%
\begin{frame}
\frametitle{\texorpdfstring{\culine[\sectiontitleulinecolor]{\sectiontitle}}{\sectiontitle} }
\framesubtitle{Pintatermi}%Surface term
\appear{}

%
\begin{itemize}
	\item \CDAlert[blue]{Pinnan nukleoneilla}\appear{ on \Alert[blue]{vähemmän puoleensavetäviä naapureita} kuin sisäosan nukleoneilla.} \appear{Joten pinnalta pitäisi olla helpompaa irrottaa nukleoni} \implies{ Pinnan nukleonien sidosenergia on alhaisempi kuin sisänukleonien }

\item Muistellaan yhtälöä
$ \quad 
\appear{V}  \equal \frac{4}{3} \appear{\pi r^{3}} \quad \Rightarrow \quad \appear{r =} \appear{\textrm{((}} \frac{3}{4 \pi} \appear{V} \appear{\textrm{))}^{1 / 3}}
$
\item[] Pinta-alaksi saadaan $\appear{=4 \pi r^{2}} \appear{=4 \pi ( \Frac{3}{4 \pi} V)^{2 / 3}} \propto$\appear{\;V$^{2 / 3}$}

\end{itemize}

% 6.5 cm = midway, 10 cm = 3/4 
\vspace{2.5mm}
\begin{minipage}{8.5cm}
\begin{itemize}
\item Sidosenergian pitäisi siis pienentyä pinta-alaan verrannollisella tekijällä\appear{, mikä on verrannollinen tilavuuteen$^{2 / 3}$}

	\item Koska \Alert[red]{tilavuus on verrannollinen nukleonien lukumäärään}\appear{, saadaan verrannollisuus:}
\[
BE \propto  \minus \appear{c_{2}} \appear{A^{2 / 3}}
\]

\end{itemize}
\end{minipage}
%\vspace{2.5mm}

\xdef\loc{south east}%
\begin{tikzpicture}[remember picture,overlay] % anchor=south
    \node (A) [yshift=-0.25*\figYshift,xshift=\figXshift, anchor=\loc] at (current page.\loc) {\includegraphics[width=5cm]{Figures_booklet/SOM_12_Nuclear_physics_figures/Tikz_SOM_Nuclear_liquid_drop_surface.pdf}};% 8cm max height
\end{tikzpicture}

\endofslide 
%\endofsection
\end{frame}
%%%%%%%%%%%%%%%

%%%%%%%%%%%%%%%
\begin{frame}
\frametitle{\texorpdfstring{\culine[\sectiontitleulinecolor]{\sectiontitle}}{\sectiontitle} }
\framesubtitle{Coulombinen termi}%Coulomb term
\appear{}

%
\begin{itemize}
	\item Coulombinen (sähköinen) vuorovaikutus protonien välillä nostaa potentiaalienergiaa  $\Rightarrow$ \appear{\Alert[red]{Coulombinen vuorovaikutus pienentää sidosenergiaa}} 

\item Ytimessä on $Z$ protonia\appear{, jotka \CDAlert[red]{hylkivät} toisia $Z-1$ protonia}

\item Coulombinen energiatermi\appear{, joka on verrannollinen termiin $\frac{e^{2}}{r}$}\appear{, vaikuttaa siis kuten alla:}
\end{itemize}

% 6.5 cm = midway, 10 cm = 3/4 
\vspace{2.5mm}
\begin{minipage}{8.5cm}
\begin{itemize}
	\item[] 
	$$
	\begin{aligned}
\frac{Z(Z-1)}{<\text{protonien etäisyys}>} &\propto \frac{Z(Z-1)}{\text{pisaran säde}}\\ 
&\propto \frac{Z(Z-1)}{A^{1 / 3}} \\
\Rightarrow \qquad  \qquad 
\appear{B E &\propto} \appear{-c_{3}} \frac{Z(Z-1)}{A^{1 / 3}} 
\end{aligned}
$$
%\item[]  $$ \Rightarrow \qquad 
%\appear{B E \propto} \appear{-c_{3}} \frac{Z(Z-1)}{A^{1 / 3}} \hspace{4cm}
%$$
\end{itemize}
\end{minipage}

\xdef\loc{south east}%
\begin{tikzpicture}[remember picture,overlay] % anchor=south
    \node (A) [yshift=-0.25*\figYshift,xshift=\figXshift, anchor=\loc] at (current page.\loc) {\includegraphics[width=5cm]{Figures_booklet/SOM_12_Nuclear_physics_figures/Tikz_SOM_Nuclear_liquid_drop_coulomb.pdf}};% 8cm max height
\end{tikzpicture}

\endofslide 
%\endofsection
\end{frame}
%%%%%%%%%%%%%%%

%%%%%%%%%%%%%%%
\begin{frame}
\frametitle{\texorpdfstring{\culine[\sectiontitleulinecolor]{\sectiontitle}}{\sectiontitle} }
\framesubtitle{Epäsymmetriatermi (Paulin termi)}%Asymmetry term
\appear{}

% 6.5 cm = midway, 10 cm = 3/4 
%\vspace{2.5mm}
\begin{minipage}{8cm}
\begin{itemize}
	\item Pienemmillä ytimillä $N \cong Z$ \appear{mutta isommilla $N>Z$.}%
	\appear{\xdef\loc{east}
\begin{tikzpicture}[remember picture,overlay] % anchor=south
    \node (A) [yshift=0*\figYshift,xshift=\figXshift, anchor=\loc] at (current page.\loc) {\includegraphics[height=8cm]{Figures_booklet/SOM_12_Nuclear_physics_figures/SOM_Nuclear_physics_Figure_13.png}}; % 8cm max height
\end{tikzpicture}\CR[ref = h]{}}
 \appear{Tämä vaikuttaa energiatasojen protonien ja neutronien väliseen miehitystasapainoon.} \appear{Ilmiö on kvanttimekaaninen ja on seurausta \CDAlert[blue]{kieltosäännöstä}}

\item Jos neutroneita on enemmän\appear{, niin suurempi määrä energiatiloja tulee miehitetyksi} \appear{(kaksi neutronia per taso)} \appear{ja systeemien ko\-ko\-nais\-e\-nergia kasvaa}

\item Ja koska meillä on prosessi joka kasvattaa energiaa\appear{, nukleonien välinen \Alert[blue]{sidosenergia pienentyy}}
\end{itemize}
\end{minipage}
%\vspace{2.5mm}

\endofslide 
%\endofsection
\end{frame}
%%%%%%%%%%%%%%%

%%%%%%%%%%%%%%%
\begin{frame}
\frametitle{\texorpdfstring{\culine[\sectiontitleulinecolor]{\sectiontitle}}{\sectiontitle} }
\framesubtitle{Epäsymmetriatermi (Paulin termi)}%Asymmetry term
\appear{}

\small
\begin{itemize}
	\item Tarkastellaan tilannetta missä neutronit ja protonit ovat miehittäneet energiatilat yhtä korkealle\appear{, $E_{F}$}\appear{, ja oletetaan energiatilat tasavälisiksi ($\delta E$)}
\appear{\xdef\loc{south east}%
\begin{tikzpicture}[remember picture,overlay]% anchor=south
    \node (A) [yshift=-0.4*\figYshift,xshift=\figXshift, anchor=\loc] at (current page.\loc) {\includegraphics[width=5cm]{Figures_booklet/SOM_12_Nuclear_physics_figures/SOM_Nuclear_physics_Figure_15.png}};% 8cm max height
\end{tikzpicture}\CR{}}%
\end{itemize}
% 6.5 cm = midway, 10 cm = 3/4 
\vspace{1.4mm}
\begin{minipage}{8.75cm}
\begin{itemize}[itemsep=1.6mm]
	\item Jos \Alert[fuchsia]{$j$ neutronia muuttuu protoneiksi} \appear{(tai päinvastoin)}\appear{, ylimpien neutroneiden täytyy siirtyä $E_{F}$:llä olevalta tilalta ylemmälle protonitilalle $E_{F}+(j / 2) \delta E$.} \appear{Muiden täytyy tehdä vastaava hyppäys.} \appear{Jos jokaisen hiukkasen energia kasvaa termillä $(j / 2) \delta E$}\appear{, kokonaisenergia kasvaa määrällä $\textrm{((}j^{2} / 2\textrm{))} \delta E$}

\item Ilmaistaan nyt $j$ \appear{neutroni- ja protonilukujen avulla ($N$ ja $Z$).} \appear{Ensin $N=Z$}\appear{, sitten tapahtuu $j$:n protonin lisäys ja $j$:n neutronin vähennys.} \appear{$Z-N$ täytyy nyt olla $2 j$}\appear{, jolloin $j=\Frac{1}{2}(Z-N)$.} \appear{Eli epäsymmetrian $Z$:n ja $N$:n välillä pitäisi pienentää sidosenergiaa (kasvattaa kokonaisenergiaa) tekijällä}
\[
(N-Z)^{2} \appear{\delta E}  \equal \appear{(A-2 Z)^{2}} \appear{\delta E} \quad \Rightarrow \quad \appear{BE} \propto \appear{-c_{4}} \frac{(A-2 Z)^{2}}{A}
\]
\appear{missä $A$:lla jakaminen normalisoi termin}
\end{itemize}
\end{minipage}
%\vspace{2.5mm}


\endofslide 
%\endofsection
\end{frame}
%%%%%%%%%%%%%%%

%%%%%%%%%%%%%%%
\begin{frame}
\frametitle{\texorpdfstring{\culine[\sectiontitleulinecolor]{\sectiontitle}}{\sectiontitle} }
\framesubtitle{Kytkeytymistermi (paritermi)}%Pairing term
\appear{}

%
\begin{itemize}
	\item Mittausten perusteella on selvää että \appear{\CDAlert[fuchsia]{vahva voima suosii} protonien ja neutronien \CDAlert[fuchsia]{kytkeytymistä}.} \appear{Tämä viimeinen termi$^{*}$ mahdollistaa paremman sovitteen tekemisen}\footnoteleft{{\appear{$^*$Huomaa että tämä 5:s (kytkeytymis)termi puuttuu Harrisin kirjasta}}}

\item Termi ottaa huomioon sen tosiasian että \Alert[fuchsia]{nukleoni–nukleoni vuorovaikutus} ei pelkästään nosta keskimääräistä nukleonien kokemaa potentiaalia\appear{, vaan että sillä on \CDAlert[blue]{samanlaisia} \CDAlert[blue]{ominaisuuksia}} \appear{kuin suprajohteissa tapahtuvissa elektroniparien välisissä vuorovaikutuksissa} \appear{(\Alert[blue]{Cooperin parien vuorovaikutus})}

\item Tämä energiatermi on:
\item[] \begin{itemize}
	\item Positiivinen \appear{(sitova)} \appear{jos sekä $Z$ että $N$ ovat parillisia,}
\item negatiivinen \appear{(hajottava)} \appear{jos sekä $Z$ että $N$ ovat parittomia}\appear{, nolla muuten}

\end{itemize}
\end{itemize}


\endofslide 
%\endofsection
\end{frame}
%%%%%%%%%%%%%%%

%%%%%%%%%%%%%%%
\begin{frame}
\frametitle{\texorpdfstring{\culine[\sectiontitleulinecolor]{\sectiontitle}}{\sectiontitle} }
\framesubtitle{Semi-empiirinen sidosenergian kaava}%Semi-empirical binding energy formula
\appear{}

\begin{itemize}
	\item Termit voidaan nyt laskea yhteen\appear{, jolloin voidaan muodostaa semi-empiirinen kaava ytimien sidosenergioille.}\appear{ Tätä kaavaa kutsutaan joskus myös  \CDAlert[fuchsia]{Weizsäckerin kaavaksi}:}
	\[
B E \equal \appear{c_{1} A}  \minus \appear{c_{2} A^{2 / 3}}  \minus \appear{c_{3} \frac{Z(Z-1)}{A^{1 / 3}}}  \minus \appear{c_{4} \frac{(A-2 Z)^{2}}{A}} \appear{\pm c_{5} A^{-4 / 3}}
\]
\item[] missä 
$$
\begin{aligned}
&\appear{c_{1}=15,8 \mathrm{\;MeV}} \\
&\appear{c_{2}=17,8 \mathrm{\;MeV}} \\
&\appear{c_{3}=0,71 \mathrm{\;MeV}} \\
&\appear{c_{4}=23,7 \mathrm{\;MeV}} \\
&\appear{c_{5}=15,8 \mathrm{\;MeV}}
\end{aligned}
$$
\end{itemize}

%\xdef\loc{south east}
%\begin{tikzpicture}[remember picture,overlay] % anchor=south
%    \node (A) [yshift=-0.25*\figYshift,xshift=\figXshift, anchor=\loc] at (current page.\loc) {\includegraphics[width=7cm]{Figures_booklet/SOM_12_Nuclear_physics_figures/SOM_Nuclear_physics_Figure_1.png}}; % 8cm max height
%\end{tikzpicture}

\endofslide 
%\endofsection
\end{frame}
%%%%%%%%%%%%%%%


%%%%%%%%%%%%%%%
\begin{frame}
\frametitle{\texorpdfstring{\culine[\sectiontitleulinecolor]{\sectiontitle}}{\sectiontitle} }
\framesubtitle{Semi-empiirinen sidosenergian kaava}%Semi-empirical binding energy formula
\appear{}

\begin{itemize}
	\item Kuva 16 osoittaa kuinka hyvin kaava \appear{(\Alert[fuchsia]{pinkki viiva})} \appear{ennustaa nukleonikohtaisen sidosenergian} \appear{eli sidososuuden (kokeelliset arvot: \Alert[black]{musta viiva})}

\item Huomaa että termeillä on fysikaalinen alkuperä\appear{ ja ne  \Alert[fuchsia]{perustuvat järkeviin oletuksiin}.}\appear{ Mutta numeeriset arvot kaavan \CDAlert[blue]{kertoimille} on \Alert[blue]{löydetty sovittamalla yhtälö mittaustuloksiin} eri nuklidien  \\sidosenergioista}
\end{itemize}

\xdef\loc{south east}
\begin{tikzpicture}[remember picture,overlay] % anchor=south
    \node (A) [yshift=-0.25*\figYshift,xshift=\figXshift, anchor=\loc] at (current page.\loc) {\includegraphics[width=7cm]{Figures_booklet/SOM_12_Nuclear_physics_figures/SOM_Nuclear_physics_Figure_16.png}}; % 8cm max height
\end{tikzpicture}
\footnote{\HarrisCRshort}

\endofslide 
%\endofsection
\end{frame}
%%%%%%%%%%%%%%%

%%%%%%%%%%%%%%%
\begin{frame}
\frametitle{\texorpdfstring{\culine[\sectiontitleulinecolor]{\sectiontitle}}{\sectiontitle} }
\framesubtitle{Kuorimalli}%Shell model
\seminormal
\begin{itemize}
	\item \appear{\xdef\loc{south east}%
\begin{tikzpicture}[remember picture,overlay] % anchor=south
    \node (A) [yshift=-0.25*\figYshift,xshift=\figXshift, anchor=\loc] at (current page.\loc) {\includegraphics[width=5cm]{Figures_booklet/SOM_12_Nuclear_physics_figures/Tikz_SOM_Nuclear_liquid_drop_just_interior.pdf}};% 8cm max height
\end{tikzpicture}}%
\CDAlert[blue]{Nestepisaramallin} lisäksi\appear{, toinen menestyksekäs malli kuvaamaan atomin ytimen rakennetta on nimeltään \CDAlert[red]{kuorimalli}}
	
	\item Huomaa \CDAlert[fuchsia]{että nämä mallit poikkeavat} toisistaan huomattavasti\appear{, mikä ilmentää vahvan voiman monimutkaisuutta}

\end{itemize}

% 6.5 cm = midway, 10 cm = 3/4 
\vspace{2.5mm}
\begin{minipage}{8.5cm}
\begin{itemize}
	\item Missä nestepisaramalli on \Alert[blue]{enimmäkseen klassinen} \appear{(paitsi epäsymmetria- ja kytkeytymistermit)}\appear{, \Alert[red]{kuorimalli pohjautuu kvanttimekaniikkaan}} 

\item Kuorimallissa hiukkasia käsitellään kvanttimekaanisina tiloina keskimääräisessä potentiaalikaivossa \appear{ja suurimmaksi osaksi unohdetaan hiukkasten väliset vuorovaikutukset}

\item Kuorimallin vahvuutena \appear{on ollut ennustaa \CDAlert[red]{ydinten kulmaliikemäärämomentit}}\appear{, ja $N$:n ja $Z$:n taipumukset koostua \CDAlert[red]{taikaluvuista}}\appear{, tai olla ainakin parillisia}
\end{itemize}
\end{minipage}
%\vspace{2.5mm}

\endofslide 
%\endofsection
\end{frame}
%%%%%%%%%%%%%%%


%%%%%%%%%%%%%%%
\begin{frame}
\frametitle{\texorpdfstring{\culine[\sectiontitleulinecolor]{\sectiontitle}}{\sectiontitle} }
\framesubtitle{Taikaluvut}%Magic numbers
%
\appear{}
% 6.5 cm = midway, 10 cm = 3/4 
%\vspace{2.5mm}
\begin{minipage}{8.5cm}
\begin{itemize}
	\item Kuorimallia ehdotettiin ratkaisuksi kun huomattiin että ydinten stabiilisuuksissa on säännönmukaisuuksia\appear{, ks.\ esim.\ kuvat joissa $N$ tai $Z$ ovat} \appear{2}\appear{, 8}\appear{, 20}\appear{, 28, 50, 82, 126}
	\appear{\xdef\loc{east}
\begin{tikzpicture}[remember picture,overlay] % anchor=south
    \node (A) [yshift=0.*\figYshift,xshift=\figXshift, anchor=\loc] at (current page.\loc) {\includegraphics[height=8cm]{Figures_booklet/SOM_12_Nuclear_physics_figures/SOM_Nuclear_physics_Figure_13.png}}; % 8cm max height
\end{tikzpicture}\CR[ref = h]{}} 
	
	\item Nämä luvut selvästikin johtavat kasvaneeseen vakauteen\appear{, kuten atomien elektronikuorten tapauksessa} \appear{2}\appear{, 10}\appear{, 18}\appear{, 36}\appear{, 54, $\ldots$}
\end{itemize}
\end{minipage}
%\vspace{2.5mm}

\endofslide 
%\endofsection
\end{frame}
%%%%%%%%%%%%%%%

%%%%%%%%%%%%%%%
\begin{frame}
\frametitle{\texorpdfstring{\culine[\sectiontitleulinecolor]{\sectiontitle}}{\sectiontitle} }
\framesubtitle{Esimerkki 3}
%
\appear{}

\begin{itemize}
	\item Virittyneessä tilassa oleva ydin voi relaksoitua emittoimalla gammasäteen \appear{(eli fotonin).} \appear{Oleta (karkeana approksimaationa) että nukleonin massaluku $A=100$}\appear{, ja se on vangittuna äärettömän syvään 1D-potentiaalikaivoon.} \appear{Kuinka suuri on emittoituneen fotonin energia (suuruusluokka-arvio riittää)?}

\item Potentiaalikaivon energiatasot ovat: \appear{$\appear{E=}\fraca{\hbar^{2} \pi^{2}}{2 m_{p} L^{2}} \appear{n^{2}}$}

\item Oletetaan $\appear{A=100} \quad \Rightarrow \quad \appear{L=2 r}\appear{=2 \times R_{0} A^{1 / 3}}\appear{=2 \times 1,2 \;fm \times 100^{1 / 3}=1,11 \times 10^{-14} \mathrm{\;m}}$ \\
\phantom{Oletetaan $\appear{A=100} \quad$}$\Rightarrow \quad \appear{r \Approx 5,5 \mathrm{\;fm}}$

\item Nukleonin massa on $\appear{m_{p}=1,67 \times 10^{-27} \mathrm{\;kg} }$

\item[] $\Rightarrow$ \appear{Energia:} $\quad \appear{E=\Frac{\hbar^{2} \pi^{2}}{2 m_{p} L^{2}} n^{2}}\appear{=3,3 \times 10^{-13} \;J \cdot n^{2}} \appear{\cong(2 \mathrm{\;MeV}) \cdot n^{2}}$

\item Huomaa että $\gamma$-fotonin energia on yleensä suuruusluokkaa $1 \mathrm{\;MeV}$\appear{, joten tulos on oikeanlainen}
\end{itemize}

%\endofslide 
\endofsection
\end{frame}
%%%%%%%%%%%%%%%

